% Section 1 — Introduction. Geosciences framing (not ML conference framing).

Ambient fine particulate matter (PM\textsubscript{2.5}) is the
single largest environmental health risk in Europe, contributing to
roughly 240{,}000 premature deaths in 2022 (European Environment
Agency, 2024). The 2024 revision of the EU Ambient Air Quality
Directive lowers the annual mean PM\textsubscript{2.5} target from
25 to 10\,$\mu$g\,m$^{-3}$, in line with the World Health Organization
guidelines, and obliges Member States to deliver dense and accurate
short-range forecasts to the public. This requirement is challenging
because PM\textsubscript{2.5} fields exhibit very strong spatial
gradients (urban hot-spots, valleys, coastlines, smoke plumes)
that are not resolved by current operational chemistry-transport
models running at 0.1$^\circ$ to 0.4$^\circ$.

The Copernicus Atmosphere Monitoring Service (CAMS) regional
ensemble \citep{cams} delivers operational forecasts at a nominal
resolution of about 10\,km. National systems such as LOTOS-EUROS
\citep{schaap2008lotos} and EMEP-MSC-W \citep{simpson2012emep} run
at comparable resolutions. While these models embed detailed
representations of atmospheric chemistry and transport, their grid
spacing remains an order of magnitude coarser than what users
require for site-level guidance, and producing a single forecast at
1\,km would multiply their compute cost by $\sim$100. At the other
end of the spectrum, satellite-derived analyses such as
GHAP \citep{wei2023ghap} provide retrospective 1\,km estimates,
but with a several-day latency and no forecast capability.

Recent deep-learning weather and air-quality models (Pangu,
\citealp{bi2023pangu}; GraphCast, \citealp{lam2023graphcast};
Aurora, \citealp{bodnar2024aurora}; ClimaX, \citealp{nguyen2023climax})
have demonstrated that data-driven approaches can rival or exceed
operational physics-based forecasts at coarse-to-medium resolution.
None of them currently produce continental fine-resolution PM
forecasts: the per-step token count of a 1\,km Vision Transformer over
Europe ($\sim 4{,}000 \times 8{,}000$ pixels) exceeds the memory of
any single accelerator. This paper introduces \textbf{CRAN-PM v1.0}
to close this gap: a multi-scale Vision Transformer that fuses
25\,km meteorology with 1\,km PM\textsubscript{2.5} analyses through
a cross-resolution attention bridge, packaged as an open-source
Python tool.

\paragraph{Contributions of this paper.}
This manuscript is a model description paper accompanying
the methodology paper \citet{kheder2026eccv}. We focus on the
software contribution and on aspects that matter for community
adoption rather than benchmark numbers.
\begin{itemize}
\item \textbf{An open-source, pip-installable Python package}
(\texttt{cranpm}) with a stable public API
(\texttt{CRANPMForecaster.from\_pretrained}, \texttt{.predict}),
a command-line interface, and a documented configuration system.
\item \textbf{Dual GPU backend support} for NVIDIA CUDA and AMD ROCm,
including container images for both, distributed training that is
portable beyond LUMI, and a reproducible benchmarking harness.
\item \textbf{Reproducibility infrastructure}: continuous integration,
pinned dependency set, deposited intermediate outputs and
seedable training/eval workflows.
\item \textbf{Three documented case studies of operational interest}:
the March 2022 Saharan dust intrusion, the July 2022 Iberian
wildfires, and the January 2022 Polish/Slovak winter heating
episode (Section~\ref{sec:cases}).
\end{itemize}

\paragraph{Scope and roadmap.}
Section~\ref{sec:related} situates CRAN-PM in the existing
software ecosystem for European air-quality forecasting.
Section~\ref{sec:model} summarises the model architecture (the
mathematical content is shared with the methodology paper).
Section~\ref{sec:software} is the heart of this manuscript: it
documents the software architecture, the GPU backend abstractions
and the reproducibility tooling. Section~\ref{sec:perf} reports
single-GPU and distributed performance benchmarks.
Section~\ref{sec:eval} extends the original evaluation with
country- and season-stratified metrics and a comparison to the
operational CAMS forecast. Section~\ref{sec:cases} walks through
the three case studies. Section~\ref{sec:ablation} reports an
ablation of the cross-attention design.
Section~\ref{sec:limits} discusses limitations.
Section~\ref{sec:availability} lists code, data, and reproduction
recipes.
